%Documentação segu# Documentação Técnica de Engenharia — Projeto Delta v4c (v1d)

Este documento descreve a arquitetura, as decisões de projeto, a modelagem matemática e os planos de validação da ferramenta computacional desenvolvida para o processamento e análise de grafos de colaboração em repositórios de código aberto (Alvo: `vuejs/core`).

---

## 1. Visão Geral da Arquitetura do Sistema
O ecossistema foi completamente unificado em Python, migrado a partir de um protótipo inicial em Java. O sistema adota a Arquitetura Orientada a Eventos (EDA) de forma assíncrona, desacoplando o pipeline de extração de dados (ETL), o motor matemático de tratamento de grafos e a camada de visualização em interface gráfica.

### Árvore Estrutural de Módulos (v1d)
* `abstract_graph.py`: Contrato base abstrato (`AbstractGraph`) e definições de enums de representação.
* `graphs.py`: Classes concretas de persistência e adjacência (`AdjacencyListGraph` e `AdjacencyMatrixGraph`).
* `orchestrator_hibrido_alpha0e.py`: Motor multithread de mineração, paginação e controle de rate-limits da API do GitHub.
* `gui_ctk.py`: Interface Gráfica de Usuário interativa desenvolvida sobre a biblioteca `CustomTkinter`.
* `data.json` / `meu_qrcode.png`: Componentes de carga de credenciais e tokens de acesso de forma criptografada.

---

## 2. Motor de Grafos e Otimização Teórica

### Classe Abstrata e API Obrigatória
A infraestrutura herda estritamente de `AbstractGraph`. Para cumprir as restrições de desempenho exigidas, a classe concreta `AdjacencyListGraph` substitui o encadeamento dinâmico comum por um arranjo indexado composto por dicionários nativos (`List[Dict[int, float]]`).

* **Complexidade Temporal:** A verificação de existência de arcos (`has_edge`) e a recuperação de pesos de arestas (`get_edge_weight`) ocorrem em tempo constante $O(1)$.
* **Restrições Estruturais:** O grafo é estritamente simples e direcionado. Não são permitidos laços (a diagonal principal é validada e zerada) nem arestas múltiplas. A inserção via `add_edge` é puramente idempotente.

### Calibração de Pesos Integrados
As conexões direcionadas ($u \rightarrow v$) possuem pesos que refletem a intensidade da colaboração técnica no ecossistema:
* **Comentário em Discussão (Issue / PR):** Peso 2.0
* **Abertura de Issue comentada por terceiros:** Peso 3.0
* **Revisão Formal ou Aprovação de Código:** Peso 4.0
* **Ato de Efetuar o Merge de um Pull Request:** Peso 5.0

---

## 3. Pipeline de Mineração e Resiliência Concorrente (EDA)

### Separação de Canais por Tripla Fila
Para mitigar condições de corrida (*race conditions*) detectadas em versões anteriores, onde workers e o loop de controle disputavam a mesma fila, o barramento `EventBus` foi refatorado utilizando três estruturas de `queue.Queue` isoladas e thread-safe:

1.  **`task_queue`**: Gerencia de forma concorrente a fila de páginas pendentes a serem mineradas pelo pool de threads.
2.  **`data_queue`**: Canal exclusivo de alta vazão para transporte de payloads de dados (`DATA_EXTRACTED`) consumidos pelo armazenador em disco.
3.  **`notification_queue`**: Linha prioritária de mensagens de infraestrutura (como alertas de `TOKEN_COOLDOWN`).

### Gerenciamento Multi-Token e Cooldown Global
O `Orchestrator` distribui as requisições HTTP síncronas entre 4 workers paralelos alimentados por um `ThreadPoolExecutor`. 
* Ao interceptar uma restrição de tráfego (HTTP 403 / 429), o `TokenManager` isola o token e calcula o tempo de liberação.
* O laço supervisor monitora a função `all_in_cooldown()`. Caso todas as chaves atinjam o teto limite, o sistema invoca `get_next\_reset\_time()` e suspende a execução de forma segura, disparando avisos na interface e impedindo o travamento do software.

---

## 4. Modelagem Matemática de Métricas Implementadas

O sistema processa a topologia da rede aplicando algoritmos estruturais nativos:

* **Centralidade de Grau ($C_D$):** Mapeia o volume de conexões diretas de entrada e saída de cada ator.
    $$C_D^{in}(u) = \sum_{v \in V} A_{vu}, \quad C_D^{out}(u) = \sum_{v \in V} A_{uv}$$
* **Centralidade de Intermediação ($C_B$):** Identifica colaboradores que atuam como pontes e elos de articulação entre subgrupos.
    $$C_B(u) = \sum_{s \neq u \neq t} \frac{\sigma_{st}(u)}{\sigma_{st}}$$
* **Centralidade de Proximidade ($C_C$):** Mede a velocidade de acesso à informação com base no inverso das distâncias mínimas (BFS).
    $$C_C(u) = \frac{|V| - 1}{\sum_{v \neq u} d(u, v)}$$
* **PageRank ($PR$):** Determina recursivamente a influência de um desenvolvedor ponderando a importância de seus vizinhos.
    $$PR(u) = \frac{1-d}{|V|} + d \sum_{v \in B_u} \frac{PR(v)}{C_D^{out}(v)}$$
* **Densidade da Rede ($\rho$):** Proporção de acoplamento técnico global da comunidade.
    $$\rho = \frac{|E|}{|V|(|V| - 1)}$$

---

## 5. Arquitetura Concorrente da Interface Gráfica (GUI)

Para evitar o congelamento da interface visual (`Main UI Thread`) durante a execução do pipeline de dados, o arquivo `gui_ctk.py` implementa um modelo assíncrono baseado em *Threading* e *Polling*:

* **UI Thread:** Executa de forma isolada o método `mainloop()` do CustomTkinter, garantindo renderização contínua a 60Hz.
* **Worker Thread:** O acionamento da mineração instancia uma thread em segundo plano (`GUI-MiningWorker`), liberando a janela para interações do usuário.
* **Mecanismo de Polling:** Através do agendador `self.after()`, o método `_poll_queues()` inspeciona a fila `log_queue` a cada 120ms, drenando os registros e capturando eventos de sistema de forma thread-safe sem travar a renderização.

---

## 6. Plano de Validação e Testes Automatizados

Para homologação das restrições de integridade da API, implementou-se uma suíte de testes unitários automatizados baseados no pacote `unittest`.

### Casos de Teste Cobertos
* `test_idempotency_add_edge`: Garante que tentativas de reinserção de uma aresta existente não gerem duplicidade nem alterem o contador `edge_count`.
* `test_loop_restriction`: Valida se o sistema lança corretamente uma exceção do tipo `ValueError` ao tentar conectar um vértice a ele mesmo.
* `test_index_out_of_bounds`: Certifica o disparo de um erro do tipo `IndexError` caso arcos façam referência a IDs fora dos limites estabelecidos do grafo.

### Como Executar os Testes Unitários
A partir da raiz do projeto no terminal, execute:
```bash
python -m unittest -v graphs.pyndo Paulo
